%-----------------------------------------
% Jose Luis Honorato
% Resumé
% Based on Michael Wright's resume
% URL: http://github.com/mdwrigh2/resume
% This document is covered under the Creative Commons Attribution 3.0 Unported License
%-----------------------------------------

%!TEX TS-program = xelatex
%!TEX encoding = UTF-8 Unicode

\documentclass[11pt,letterpaper,final]{moderncv}
\usepackage{fontspec}
\usepackage{wrapfig}
\usepackage{graphics}
%\usepackage{graphicx}
\moderncvtheme{classic}

% DOCUMENT LAYOUT
\usepackage[scale=0.85]{geometry}
\setlength{\hintscolumnwidth}{2cm}
\AtBeginDocument{\recomputelengths}
% FONTS
% \usepackage[utf8]{inputenc}
\defaultfontfeatures{Mapping=tex-text} % converts LaTeX specials (``quotes'' --- dashes etc.) to unicode

% Remove % to set to different fonts
%\setromanfont [Ligatures={Common},Numbers={OldStyle}]{Adobe Caslon Pro}
%\setmonofont[Scale=0.8]{Monaco} 

% ---- MARGIN años
%\newcommand{\años}[1]{\marginpar{\scriptsize #1}}


% PDF SETUP
% ---- FILL IN HERE THE DOC TITLE AND AUTHOR


% Personal Information 
\firstname{Jose L.}
\familyname{Honorato}
\address{4383 Pdte. Errazuriz, dep 902}{Las Condes, Santiago, Chile}
\mobile{+56 9 92619739}
\email{jlhonora@ing.puc.cl}
%\homepage{www.derefnull.com/about}

\title{Ingeniero Civil Eléctrico, M. Sc.}

\nopagenumbers{}

\begin{document}

\maketitle 

\section{Estudios}
\cventry{Junio 2010 - Agosto 2011}{Magíster en Ciencias de la Ingeniería, PUC}{Santiago, Chile}{}{}{}%{Graduado con \textit{distinción máxima}}
\cventry{Marzo 2005 - Junio 2010}{Ingeniero Civil Electricista, PUC}{Santiago, Chile}{}{}{}
\section{Experiencia} 
  \cventry{Julio 2011 - Presente}{Ingeniero I+D - Co-fundador}{Satelinx}{Santiago, Chile}{}{
    \begin{itemize}
      \item Diseño de hardware, enfocado en equipos GPS y de comunicación satelital Iridium. 
	  \item Responsable de la programación de firmware de los equipos, el diseño de PCB y testing automático de software y hardware. Trabajo mano a mano con el equipo de manufactura. 
    \end{itemize}
  }
  \cventry{Agosto 2011 - Presente}{Profesor part-time}{Departamento de Computación, PUC}{Santiago, Chile}{}{
    \begin{itemize}
      \item Profesor del curso \"Laboratorio de Arquitectura de Computadores\". Aprendizaje de software embebido, con plataformas como msp430, Raspberry Pi, BeagleBoard. Responsable del re-diseño del curso para el uso de tecnologías más actuales.  
    \end{itemize}
  }
  \cventry{Junio 2010 - Agosto 2011}{Investigador}{Centro de Imágenes Biomédicas}{Santiago, Chile}{}{
    \begin{itemize}
      \item Diseño de un algoritmo de procesamiento de imágenes para generar mejores diagnósticos en imágenes con grasa y agua. 
      \item Trabajo publicado en la mejor revista internacional de imágenes resonancia magnética.
    \end{itemize}
  }
  \cventry{Verano 2010}{Software Engineering Intern}{Infosys}{Bangalore, India}{}{
    \begin{itemize}
      \item Desarrollo de un plan de negocios relacionado a la publicidad en dispositivos móviles. 
    \end{itemize}
  }
\section{Otros proyectos}
  \cvlistitem{Co-fundador Mobile and Embedded Computing Laboratory. Primera institución chilena aceptada en el iPhone Developer University Program.}
  \cvlistitem{Participante de The Sika Solar Experience. Invitado por Sika A.G. a visitas a plantas fotovoltaicas alrededor de Europa.}
\section{Habilidades profesionales} 
  \cvline{Lenguajes:}{C (2 años), Matlab (4 años), Python (2 años), C\# (4 años), Bash (1 año)}
  \cvline{Hardware:}{PCB layout, diseño RF. Protocolos: SPI, I2C, USB, UART.}
  \cvline{Plataformas:}{Mac OS X (unix) y Ubuntu Linux}
  \cvline{Misc:}{Procesamiento de imágenes, integración de sensores.}
\section{Premios}
  \cvlistitem{Mejor Ingeniero Eléctrico, \textit{Colegio de Ingenieros}, Junio 2012}
  \cvlistitem{Magna Cum Laude Award, \textit{ISMRM}, Mayo 2012}
  \cvlistitem{Mejor Tesis de Ingeniería Eléctrica, \textit{Depto. Ingeniería Eléctrica}, Sept. 2011}
  \cvlistitem{Best Intern, \textit{Infosys}, Marzo 2010}
\section{Extras}
  \cvlistitem{LinkedIn: \href{http://www.linkedin.com/in/jlhonora}{http://www.linkedin.com/in/jlhonora}}
  \cvlistitem{Github: \href{http://github.com/jlhonora}{http://github.com/jlhonora}}
  \cvlistitem{Satelinx: \href{http://satelinx.com}{http://satelinx.com}}
  \cvlistitem{Copia de mi CV: \href{http://github.com/jlhonora/resume}{http://github.com/jlhonora/resume}}
  \begin{wrapfigure}{r}{0.25\textwidth}
     \vspace{-3.3cm}
     \begin{center}
     Escanea el código\\
     para tener mi contacto
       \includegraphics[width=0.25\textwidth]{qr}
     \end{center}
  \end{wrapfigure}

\end{document}
